\documentclass[aspectratio=169,10pt,english]{beamer}

\input{../../../_preambulo_beamer.tex}
\input{../../../_comandos_beamer.tex}

\selectlanguage{english}

\title{MA1001B - Statistical Modeling for Decision Making}
\subtitle{Lesson 48: Regression Inference, Residuals, and Decisions}
\author{}
\institute{Tecnol\'ogico de Monterrey}
\date{}

\begin{document}

\begin{frame}[plain]
  \vspace{1.2cm}
  {\Large\bfseries MA1001B - Statistical Modeling for Decision Making\par}
  \vspace{0.7cm}
  {\large Lesson 48: Regression Inference, Residuals, and Decisions\par}
  \vspace{0.7cm}
  {\color{TecRojo}\rule{\textwidth}{1.2pt}}
  \vspace{0.8cm}

  MA1001B Faculty\\[0.3cm]
  Term\\[0.3cm]
  Tecnol\'ogico de Monterrey
\end{frame}

\begin{frame}{The Inference Question}
  \begin{alertblock}{Guiding question}
    Is a fitted overtime slope different from 0, and can one emergency
    weekend create that significance by itself?
  \end{alertblock}

  \vspace{0.6em}
  By the end of this lesson, you can:
  \begin{itemize}
    \item compute $SE(b_1)$ and the $t$ test of $H_0:\beta_1=0$;
    \item read a residual-versus-fitted plot;
    \item refit after deleting an influential week;
    \item refuse to treat one weekend as ordinary-week policy.
  \end{itemize}

  \vfill
  \scriptsize All overtime hours and unit counts used today are synthetic. No real company data are used.
\end{frame}

\begin{frame}{Thirty-Nine Ordinary Weeks and One Emergency}
  \small
  \begin{center}
    \begin{tabular}{lr}
      \toprule
      \textbf{Piece} & \textbf{In this sample} \\
      \midrule
      Ordinary weeks & 39, overtime in $[6,\ 22]$ \\
      Emergency weekend & $(x,y)=(48,\ 42)$ \\
      $H_0$ & $\beta_1=0$ \\
      $H_1$ & $\beta_1\neq 0$ \\
      \bottomrule
    \end{tabular}
  \end{center}

  \vspace{0.5em}
  \begin{exampleblock}{Why the weekend matters}
    A high-leverage $x$ can tilt the least-squares line. Inference on all 40
    weeks is not automatically inference about ordinary operations.
  \end{exampleblock}
\end{frame}

\begin{frame}[fragile]{Step 1: Slope, SE, and $t$}
  \begin{columns}[T]
    \begin{column}{0.40\textwidth}
      \small
      \begin{lstlisting}
t = slope / se_slope
p = two_sided_t(t, df=38)
      \end{lstlisting}

      \begin{block}{Verified output}
        $b_1=0.434381$\\
        $SE=0.072494$, $t=5.991932$\\
        $p=5.839701\times10^{-7}$
      \end{block}
    \end{column}
    \begin{column}{0.60\textwidth}
      \centering
      \includegraphics[width=\textwidth]{../figures/regression_inference_01_slope_se_t.png}
      \vspace{0.2em}
      \scriptsize $n=40$ fitted line from Step 1
    \end{column}
  \end{columns}
\end{frame}

\begin{frame}{Residuals Locate the Unusual Week}
  \small
  Residual mean is $0.000000$ by the least-squares identity. Residual SD is
  $3.140288$. The emergency-weekend residual is $8.570594$, the largest
  absolute residual in the sample.
  \[
    e_i=y_i-\hat y_i=y_i-(12.579109+0.434381\,x_i).
  \]
  A residual plot is a decision tool: it can show a fan, a curve, or one week
  that does not belong with the others.
\end{frame}

\begin{frame}[fragile]{Step 2: Residual versus Fitted}
  \begin{columns}[T]
    \begin{column}{0.40\textwidth}
      \small
      \begin{lstlisting}
e = y - (b0 + b1 * x)
      \end{lstlisting}

      \begin{block}{Verified output}
        Residual mean $0.000000$\\
        Emergency $e=8.570594$\\
        Largest $|e|$ is that week
      \end{block}
    \end{column}
    \begin{column}{0.60\textwidth}
      \centering
      \includegraphics[width=\textwidth]{../figures/regression_inference_02_residual_plot.png}
      \vspace{0.2em}
      \scriptsize Residual plot from Step 2
    \end{column}
  \end{columns}
\end{frame}

\begin{frame}{Step 3: The Significant Slope Vanishes}
  \begin{columns}[T]
    \begin{column}{0.42\textwidth}
      \small
      Without the weekend: $n=39$, $b_1=0.088445$, $t=1.070131$,
      $p=0.291492$, $R^2=0.030022$.

      \vspace{0.4em}
      Do not reject $H_0:\beta_1=0$.

      \vspace{0.5em}
      \textbf{Limit:} an influential point can create a significant slope
      that vanishes when removed. The $n=40$ test is not a decision about
      ordinary overtime weeks.
    \end{column}
    \begin{column}{0.58\textwidth}
      \centering
      \includegraphics[width=\textwidth]{../figures/regression_inference_03_influential_point.png}
      \vspace{0.2em}
      \scriptsize $n=40$ versus $n=39$ lines from Step 3
    \end{column}
  \end{columns}
\end{frame}

\begin{frame}{Concept Check}
  \begin{enumerate}
    \item Write the $t$ statistic for $H_0:\beta_1=0$ in terms of $b_1$ and
      $SE(b_1)$.
    \item Why is residual mean 0 not a diagnostic pass by itself?
    \item What happens to $p$ when the emergency weekend is deleted?
    \item Why is the $n=40$ slope the wrong ordinary-week policy number?
  \end{enumerate}

  \vspace{0.6em}
  \begin{block}{Connection}
    This closes the simple-regression sequence. Multiple predictors and
    designed experiments are later courses; the limit here is influence.
  \end{block}

  \vfill
  \scriptsize Source: OpenStax, \textit{Introductory Business Statistics 2e},
  Sections 13.2 and 13.6. CC BY-NC-SA.
\end{frame}

\appendix



\end{document}
